\section*{Chapter \ref{ch:g11:deriv} --- Differentiation: A First Course}

\begin{solution}{exo:g11:deriv:1}
For $f(x) = x^2$ at $a = 2$:
\[
\frac{(2+h)^2 - 4}{h} = \frac{4h + h^2}{h} = 4 + h
\xrightarrow[h\to0]{} 4,
\]
so $f'(2) = 4$. For $g(x) = x^2 + 3x$ at $a = 1$: $g(1) = 4$ and
\[
\frac{(1+h)^2 + 3(1+h) - 4}{h} = \frac{5h + h^2}{h} = 5 + h
\xrightarrow[h\to0]{} 5,
\]
so $g'(1) = 5$.
\end{solution}

\begin{solution}{exo:g11:deriv:2}
$f'(x) = 12x^2 - 4x + 1$.

$g'(x) = \frac{5}{2\sqrt x} - \frac{3}{x^2}$ (on $\intoo{0}{+\infty}$).

$h'(x) = 2(x^2 - 3) + (2x+1)(2x) = 2x^2 - 6 + 4x^2 + 2x
= 6x^2 + 2x - 6$ (product rule).
\end{solution}

\begin{solution}{exo:g11:deriv:3}
$f(x) = x^2 - 3x$, $a = 2$: $f(2) = -2$, $f'(x) = 2x - 3$ so $f'(2) = 1$;
\omterm{def:g11:deriv:tangent}{tangent} $y = -2 + (x - 2) = x - 4$.

$f(x) = \frac1x$, $a = 1$: $f(1) = 1$, $f'(1) = -1$; \omterm{def:g11:deriv:tangent}{tangent}
$y = 1 - (x - 1) = 2 - x$.

$f(x) = \sqrt x$, $a = 4$: $f(4) = 2$, $f'(4) = \frac{1}{2\sqrt4} =
\frac14$; \omterm{def:g11:deriv:tangent}{tangent} $y = 2 + \frac14(x - 4) = \frac{x}{4} + 1$.
\end{solution}

\begin{solution}{exo:g11:deriv:4}
$f'(x) = \frac{(x-1) - (x+2)}{(x-1)^2} = \frac{-3}{(x-1)^2}$.

$g'(x) = \frac{2x(x^2+1) - x^2 \times 2x}{(x^2+1)^2}
= \frac{2x}{(x^2+1)^2}$.

$h'(x) = -\frac{2x+1}{(x^2+x+1)^2}$ (quotient rule with $u = 1$).
\end{solution}

\begin{solution}{exo:g11:deriv:5}
$f'(x) = 3x^2 - 12x + 9 = 3(x^2 - 4x + 3) = 3(x-1)(x-3)$: positive
outside $\intcc{1}{3}$, negative inside. So $f$ is \omterm{def:g11:func:monotone}{increasing} on
$\intoc{-\infty}{1}$, \omterm{def:g10:functions:variations}{decreasing} on $\intcc{1}{3}$, \omterm{def:g11:func:monotone}{increasing} on
$\intco{3}{+\infty}$, with a local \omterm{def:g10:functions:extrema}{maximum} $f(1) = 1 - 6 + 9 - 2 = 2$ and
a local \omterm{def:g10:functions:extrema}{minimum} $f(3) = 27 - 54 + 27 - 2 = -2$.
\end{solution}

\begin{solution}{exo:g11:deriv:6}
On $\intoo{-1}{+\infty}$:
\[
f'(x) = \frac{2x(x+1) - (x^2+3)}{(x+1)^2}
= \frac{x^2 + 2x - 3}{(x+1)^2}
= \frac{(x+3)(x-1)}{(x+1)^2}.
\]
For $x > -1$, both $x + 3$ and $(x+1)^2$ are positive, so $f'(x)$ has the
sign of $x - 1$: $f$ is \omterm{def:g10:functions:variations}{decreasing} on $\intoc{-1}{1}$ and \omterm{def:g11:func:monotone}{increasing} on
$\intco{1}{+\infty}$, with \omterm{def:g10:functions:extrema}{minimum} $f(1) = \frac{4}{2} = 2$.
\end{solution}

\begin{solution}{exo:g11:deriv:7}
$f'(x) = 2x - 4$.

\emph{1.} Horizontal \omterm{def:g11:deriv:tangent}{tangent}: $f'(x) = 0$ at $x = 2$; the point is
$(2, f(2)) = (2, 1)$ (the vertex of the \omterm{def:g11:quad:parabola}{parabola}).

\emph{2.} Parallel to $y = 2x + 1$ means \omterm{def:g10:reffunc:affine}{slope} $2$: $2x - 4 = 2$, so
$x = 3$, giving the point $(3, f(3)) = (3, 2)$.
\end{solution}

\begin{solution}{exo:g11:deriv:8}
With width $x$ (two fenced sides) and length $60 - 2x$:
$A(x) = x(60 - 2x) = 60x - 2x^2$ on $\intoo{0}{30}$. Then
$A'(x) = 60 - 4x$, positive before $x = 15$ and negative after: the area
is maximal for $x = 15$, giving a $15 \times 30$ m enclosure of area
$450$ m$^2$. The vertex formula $\alpha = -\frac{60}{2 \times (-2)} = 15$
of \cref{ch:g11:quad} gives the same answer without calculus.
\end{solution}

\begin{solution}{exo:g11:deriv:9}
Let $f(x) = \frac{x+1}{2} - \sqrt x$ on $\intoo{0}{+\infty}$. Then
\[
f'(x) = \frac12 - \frac{1}{2\sqrt x} = \frac{\sqrt x - 1}{2\sqrt x},
\]
which has the sign of $\sqrt x - 1$: negative on $\intoo{0}{1}$, positive
on $\intoo{1}{+\infty}$. So $f$ decreases then increases, with \omterm{def:g10:functions:extrema}{minimum}
$f(1) = 1 - 1 = 0$. Hence $f(x) \geq 0$ for all $x > 0$, that is
$\sqrt x \leq \frac{x+1}{2}$, with equality only at $x = 1$.
\end{solution}

\begin{solution}{exo:g11:deriv:10}
From $\pi r^2 h = 500$, $h = \frac{500}{\pi r^2}$, so
\[
S(r) = 2\pi r^2 + 2\pi r \cdot \frac{500}{\pi r^2}
= 2\pi r^2 + \frac{1000}{r},
\qquad
S'(r) = 4\pi r - \frac{1000}{r^2}.
\]
$S'(r) = 0$ when $4\pi r^3 = 1000$, i.e.\ $r^3 = \frac{250}{\pi}$, so
$r = \sqrt[3]{250/\pi}$ ($\approx 4.3$ cm). Since $S'(r) =
\frac{4\pi r^3 - 1000}{r^2}$ is negative before this value and positive
after, it is a \omterm{def:g10:functions:extrema}{minimum}.
\end{solution}

\begin{solution}{exo:g11:deriv:11}
The \omterm{def:g11:deriv:tangent}{tangent} to $y = x^2$ at the point of \omterm{def:g10:coordgeom:system}{abscissa} $a$ is
$y = a^2 + 2a(x - a) = 2ax - a^2$. It passes through $P(1, -3)$ when
$-3 = 2a - a^2$, i.e.\ $a^2 - 2a - 3 = 0$, whose \omterm{def:g11:quad:discriminant}{roots} are $a = 3$ and
$a = -1$ ($\Delta = 16$). There are therefore \emph{two} \omterm{def:g11:deriv:tangent}{tangents} through
$P$:
\[
y = 6x - 9 \quad (a = 3),
\qquad
y = -2x - 1 \quad (a = -1).
\]
\end{solution}

\begin{solution}{exo:g11:deriv:12}
From \cref{ex:g11:deriv:study}, $f$ increases up to the local \omterm{def:g10:functions:extrema}{maximum}
$f(-1) = 3$, decreases to the local \omterm{def:g10:functions:extrema}{minimum} $f(1) = -1$, then increases
again; for $x$ large positive (resp.\ negative), $x^3$ dominates and $f$
takes arbitrarily large (resp.\ small) values. Reading horizontal lines
$y = k$ across the \omterm{def:g10:functions:table}{variation table}:
\begin{itemize}
  \item $k > 3$ or $k < -1$: the line meets the curve once
        --- $1$ solution;
  \item $k = 3$ or $k = -1$: the line touches a turning point and crosses
        one other branch --- $2$ solutions;
  \item $-1 < k < 3$: the line crosses all three branches --- $3$
        solutions.
\end{itemize}
\end{solution}

\begin{solution}{pb:g11:deriv:1}
\textbf{1.} $f'(x) = 6x - 5$; $g'(x) = 3x^2 - 12$. Definition:
$\frac{(1 + h)^2 - 1}{h} = 2 + h \to 2$: the \omterm{def:g11:deriv:derivative}{derivative} of
$x^2$ at $1$ is $2$.

\textbf{2.} $g(2) = 8 - 24 = -16$ and $g'(2) = 12 - 12 = 0$:
the \omterm{def:g11:deriv:tangent}{tangent} is the horizontal line $y = -16$. A horizontal
\omterm{def:g11:deriv:tangent}{tangent} signals a critical point --- here a local \omterm{def:g10:functions:extrema}{minimum}, as
the \omterm{def:g10:functions:table}{variation table} of question 5 confirms.

\textbf{3.} \omterm{def:g10:reffunc:affine}{Slope} $2a$ at $(a, a^2)$:
$y = a^2 + 2a(x - a) = 2ax - a^2$.

\textbf{4.} At $x = 0$: $y = -a^2$: the \omterm{def:g11:deriv:tangent}{tangent} crosses the
axis at depth $a^2$ below the origin --- exactly as far
\emph{below} zero as $P$ sits above it. Recipe: to draw the
\omterm{def:g11:deriv:tangent}{tangent} at a point of height $h$, join it to the axis point at
depth $h$ below the origin. (No \omterm{def:g11:deriv:derivative}{derivative} needed at the
drawing board.)

\textbf{5.} $g'(x) = 3(x^2 - 4)$: positive on
$\intoo{-\infty}{-2}$, negative on $\intoo{-2}{2}$, positive
beyond. $g$ increases to a local \omterm{def:g10:functions:extrema}{maximum} $g(-2) = 16$, falls
to a local \omterm{def:g10:functions:extrema}{minimum} $g(2) = -16$, then increases.

\textbf{6.} $F\left(0, \frac14\right)$ and $T(0, -a^2)$:
$FT = \frac14 + a^2$.

\textbf{7.} $FP = a^2 + \frac14 = FT$: isosceles at $F$, so the
base angles are equal: $\widehat{FTP} = \widehat{FPT}$.

\textbf{8.} The incoming ray at $P$ is vertical, hence parallel
to the line $(TF)$ (the $y$-axis). The \omterm{def:g11:deriv:tangent}{tangent} $(TP)$ crosses
both parallels, so the angle between the ray and the \omterm{def:g11:deriv:tangent}{tangent}
equals $\widehat{FTP}$ (alternate angles) --- which equals
$\widehat{FPT}$, the angle between $[PF]$ and the \omterm{def:g11:deriv:tangent}{tangent}.
Equal angles on the two sides of the \omterm{def:g11:deriv:tangent}{tangent} at $P$.

\textbf{9.} By the law of reflection, the ray arriving
vertically at $P$ leaves along the line making the equal angle
on the other side --- by question 8, along $[PF]$: through the
\omterm{pb:g11:quad:1}{focus}, whatever $a$ is. Reversed, light born at $F$ leaves
every point of the mirror parallel to the axis. Dishes collect,
headlights beam: proved.

\textbf{10.} $a = 1$: \omterm{def:g11:deriv:tangent}{tangent} $y = 2x - 1$, $T(0, -1)$;
$FP = \sqrt{1 + \left(1 - \frac14\right)^2} =
\sqrt{\frac{25}{16}} = \frac54$; $FT = \frac14 + 1 = \frac54$:
isosceles, as promised.

\textbf{11.} The \omterm{def:g11:deriv:tangent}{tangent} at $x_n$ is
$y = f(x_n) + f'(x_n)(x - x_n)$; it crosses $y = 0$ where
$x = x_n - \frac{f(x_n)}{f'(x_n)}$.

\textbf{12.} $x_1 = 1 - \frac{-1}{2} = \frac32$;
$x_2 = \frac32 - \frac{1/4}{3} = \frac{17}{12}$;
$x_3 = \frac{577}{408}$: Heron's \omterm{def:g10:numbers:approx}{approximations} of $\sqrt2$,
i.e.\ the fixed-point machine of \cref{pb:g10:functions:1}.
Identity: $x - \frac{x^2 - 2}{2x} = \frac{2x^2 - x^2 + 2}{2x}
= \frac{x^2 + 2}{2x} = \frac12\left(x + \frac2x\right)$ ---
Heron's recipe \emph{is} Newton's method applied to
$x^2 - 2$, two millennia early.

\textbf{13.} $x_1 = 5 - \frac{125 - 100}{75} = \frac{14}{3}
\approx 4.6667$; $x_2 \approx 4.6667 - \frac{1.63}{65.33}
\approx 4.6417$. Calculator: $\sqrt[3]{100} = 4.6416\ldots$ ---
four correct decimals in two steps.

\textbf{14.} $f'(2) = 0$: the \omterm{def:g11:deriv:tangent}{tangent} at the starting point is
horizontal (question 2) and never crosses the axis --- the
formula divides by zero. Caveat: start Newton away from
critical points (and, more generally, close enough to the \omterm{def:g11:quad:discriminant}{root}
you want).

\textbf{15.} Bisection gains one binary digit per step; Newton
roughly \emph{doubles} the number of correct digits each step
($1 \to 3 \to 6$ here). When each step is costly --- millions
of parameters, real-time navigation --- doubling beats
inching, which is why the \omterm{def:g11:deriv:tangent}{tangent} line runs inside so much
silicon.

\textbf{16.} $d^2 = 900 - w^2$, so
$S(w) = w(900 - w^2) = 900w - w^3$, for $0 < w < 30$.

\textbf{17.} $S'(w) = 900 - 3w^2 = 0$ at
$w = \sqrt{300} = 10\sqrt3 \approx 17.3$ cm ($S' > 0$ before,
$< 0$ after: a \omterm{def:g10:functions:extrema}{maximum}). Then
$d = \sqrt{600} = 10\sqrt6 \approx 24.5$ cm.

\textbf{18.} $\frac{d^2}{w^2} = \frac{600}{300} = 2$, so
$d = w\sqrt2$: depth is width times $\sqrt2$ --- the
carpenter's ratio (and the diameter-in-thirds construction
delivers it with a compass, no algebra on site).

\textbf{19.} Square beam: $w = d = \sqrt{450} \approx 21.2$
cm, stiffness $w^3 = 450^{3/2} \approx 9\,546$. Optimal beam:
$S = 10\sqrt3 \times 600 = 6\,000\sqrt3 \approx 10\,392$. Gain:
about $8.9\,\%$ stiffer from the same log --- the \omterm{def:g11:deriv:derivative}{derivative}
pays for the sawmill.

\textbf{20.} Geometry: the \omterm{def:g11:deriv:tangent}{tangent}'s equal angles turned a
paraboloid into a light-gatherer. Numerics: sliding down
\omterm{def:g11:deriv:tangent}{tangents} solves \omterm{def:g10:algebra:equation}{equations} at digit-doubling speed.
Optimization: a vanishing \omterm{def:g11:deriv:derivative}{derivative} found the log's strongest
rectangle. One object, three trades --- and grade 12's
convexity will tell, in addition, on which side of every
\omterm{def:g11:deriv:tangent}{tangent} the curve lies.
\end{solution}
